\section{Pattern-forming condition for $K$ populations}
\label{sec:K-pattern-condition}

We now state the connectivity conditions under which a $K$-population
Gaussian neural field can undergo the finite-wavenumber Hopf bifurcation
assumed in \eqref{eq:Hopf_assumption}.  Set $d=K$ and let
\begin{equation}
 \mathsf W(p)=\mathsf W_0+p\mathsf W_1
 \in\mathbb R^{K\times K}
 \label{eq:K-strength-matrix}
\end{equation}
be the signed integrated-strength matrix.  Rows are target populations and
columns are source populations.  Let
\begin{equation}
 \Sigma=(\Sigma_{ij})\in[0,\infty)^{K\times K}
 \label{eq:K-variance-matrix}
\end{equation}
be the Gaussian-variance matrix.  Thus the Fourier connectivity at spatial
mode $h$ is
\begin{equation}
 \widehat{\mathsf J}(h,p)
 =\mathsf W(p)\odot
 \exp\left(-\frac{h^2}{2}\Sigma\right),
 \qquad
 \widehat{\mathsf J}_{ij}(h,p)
 =W_{ij}(p)\exp\left(-\frac{\Sigma_{ij}h^2}{2}\right).
 \label{eq:K-Fourier-connectivity}
\end{equation}
Here the exponential and product are entrywise.  If $\sigma_{ij}$ denotes
the standard deviation instead, then $\Sigma_{ij}=\sigma_{ij}^2$.

Let
\begin{equation}
 \mathsf T=\operatorname{diag}(\tau_1,\ldots,\tau_K),
 \qquad
 \mathsf D_*=\operatorname{diag}(g_1,\ldots,g_K),
 \qquad g_j>0,
 \label{eq:K-time-gain-matrices}
\end{equation}
where $g_j$ is the derivative of the $j$th firing-rate function at the
homogeneous equilibrium.  For the convention in which each firing rate is
applied before convolution, the linear Fourier matrix is
\begin{equation}
 M(h,p)=\mathsf T^{-1}
 \left[-I+\widehat{\mathsf J}(h,p)\mathsf D_*\right].
 \label{eq:K-linear-matrix}
\end{equation}
If the firing rate is applied after recurrent summation, replace
$\widehat{\mathsf J}(h,p)\mathsf D_*$ by
$\mathsf D_*\widehat{\mathsf J}(h,p)$.  The conditions below are otherwise
unchanged.

Define the spectral growth rate
\begin{equation}
 s(h,p)=max_{\lambda\in\operatorname{spec}M(h,p)}
 \operatorname{Re}\lambda.
 \label{eq:K-spectral-growth}
\end{equation}
The connectivity is capable of producing a spatial pattern if, for some
parameter value, the homogeneous mode is stable while a nonzero spatial mode
is unstable:
\begin{equation}
 s(0,p)<0,
 \qquad
 s(h,p)>0
 \quad\text{for some }h\neq0.
 \label{eq:K-pattern-capability}
\end{equation}
At the first pattern-forming threshold $p=p_c$, there must be a
$k_0\neq0$ such that
\begin{equation}
 s(0,p_c)<0,
 \qquad
 s(k_0,p_c)=0,
 \qquad
 s(h,p_c)<0
 \quad\text{for }h\notin\{k_0,-k_0\}.
 \label{eq:K-first-pattern-threshold}
\end{equation}
After replacing $p$ by $p-p_c$, the critical value is $p=0$, as assumed in
the preceding sections.

\subsection{Coefficient condition for a finite-wavenumber Hopf bifurcation}

For every spatial mode, write
\begin{equation}
 P(\lambda;h,p)
 =\det\left(\lambda I-M(h,p)\right)
 =\lambda^K+\gamma_1(h,p)\lambda^{K-1}
 +\cdots+\gamma_K(h,p).
 \label{eq:K-characteristic-polynomial}
\end{equation}
The coefficients are determined directly by $\mathsf W(p)$ and $\Sigma$:
\begin{equation}
 \gamma_r(h,p)
 =(-1)^r
 \sum_{\substack{I\subset\{1,\ldots,K\}\\|I|=r}}
 \det M(h,p)[I,I],
 \qquad 1\leq r\leq K,
 \label{eq:K-characteristic-coefficients}
\end{equation}
where $M[I,I]$ denotes a principal submatrix.

Set $\gamma_0=1$ and $\gamma_r=0$ for $r<0$ or $r>K$.  Define the Hurwitz
determinants
\begin{equation}
 \Delta_r(h,p)
 =\det\left[\gamma_{2j-i}(h,p)\right]_{i,j=1}^{r},
 \qquad 1\leq r\leq K.
 \label{eq:K-Hurwitz-determinants}
\end{equation}
For a fixed spatial mode $h$, all eigenvalues are stable if and only if
\begin{equation}
 \Delta_r(h,p)>0,
 \qquad r=1,\ldots,K.
 \label{eq:K-Hurwitz-stability}
\end{equation}

At a finite-wavenumber Hopf point, the critical polynomial has the exact
factorization
\begin{equation}
 P(\lambda;k_0,p_c)
 =(\lambda^2+\omega_0^2)
 \left(
 \lambda^{K-2}+\beta_1\lambda^{K-3}
 +\cdots+\beta_{K-2}
 \right),
 \qquad \omega_0>0,
 \label{eq:K-Hopf-factorization}
\end{equation}
and the degree-$(K-2)$ factor must be Hurwitz.  This factorization is the
necessary and sufficient coefficient condition for one critical pair
$\pm\ii\omega_0$ and all remaining eigenvalues in the open left half-plane.

Equivalently, the critical frequency and coefficients satisfy
\begin{equation}
 \sum_{\substack{0\leq r\leq K\\K-r\ \mathrm{even}}}
 (-1)^{(K-r)/2}\gamma_r(k_0,p_c)\omega_0^{K-r}=0
 \label{eq:K-Hopf-real}
\end{equation}
and
\begin{equation}
 \sum_{\substack{0\leq r\leq K\\K-r\ \mathrm{odd}}}
 (-1)^{(K-r-1)/2}\gamma_r(k_0,p_c)\omega_0^{K-r}=0.
 \label{eq:K-Hopf-imaginary}
\end{equation}
At a generic Hopf boundary reached from the stable side, these conditions can
be detected by
\begin{equation}
 \begin{aligned}
  \Delta_r(k_0,p_c)&>0
  &&\text{for }1\leq r\leq K-2,\\
  \Delta_{K-1}(k_0,p_c)&=0,
  &\qquad \gamma_K(k_0,p_c)&>0,
 \end{aligned}
 \label{eq:K-Hopf-Hurwitz-boundary}
\end{equation}
provided no other eigenvalue is critical.  The complete pattern-forming
condition is obtained by combining \eqref{eq:K-Hopf-factorization} with
\begin{equation}
 \Delta_r(0,p_c)>0,
 \qquad
 \Delta_r(h,p_c)>0
 \quad\text{for }h\notin\{k_0,-k_0\},
 \qquad
 r=1,\ldots,K.
 \label{eq:K-noncritical-mode-stability}
\end{equation}
The bifurcation is transversal when
\begin{equation}
 \operatorname{Re}\left[
 \Psi^*M_p(k_0,p_c)\Phi
 \right]\neq0,
 \label{eq:K-pattern-transversality}
\end{equation}
where
\begin{equation}
 M(k_0,p_c)\Phi=\ii\omega_0\Phi,
 \qquad
 M(k_0,p_c)^*\Psi=-\ii\omega_0\Psi,
 \qquad
 \Psi^*\Phi=1.
 \label{eq:K-pattern-eigenvectors}
\end{equation}

For a continuous wavenumber, an interior critical mode also satisfies
\begin{equation}
 \operatorname{Re}\left[
 \Psi^*M_h(k_0,p_c)\Phi
 \right]=0,
 \qquad
 \partial_h^2\operatorname{Re}\lambda(k_0,p_c)<0.
 \label{eq:K-wavenumber-maximum}
\end{equation}
The width--strength dependence in the first equation is explicit:
\begin{equation}
 M_h(h,p)
 =-h\mathsf T^{-1}
 \left[
 \left(\mathsf W(p)\odot\Sigma\right)
 \odot\exp\left(-\frac{h^2}{2}\Sigma\right)
 \right]\mathsf D_*.
 \label{eq:K-wavenumber-derivative}
\end{equation}
On the periodic domain, the modes are discrete, so
\eqref{eq:K-first-pattern-threshold} replaces the derivative condition.

\subsection{Relation to narrower excitation}

Let $\mathcal E$ and $\mathcal I$ be the excitatory and inhibitory source
populations.  Dale's sign condition is
\begin{equation}
 W_{ij}(p)\geq0
 \quad\text{for }j\in\mathcal E,
 \qquad
 W_{ij}(p)\leq0
 \quad\text{for }j\in\mathcal I.
 \label{eq:K-Dale-condition}
\end{equation}
A direct $K$-population analogue of narrower excitation is
\begin{equation}
 \max_{j\in\mathcal E}\Sigma_{ij}
 <
 \min_{j\in\mathcal I}\Sigma_{ij},
 \qquad i=1,\ldots,K.
 \label{eq:K-narrower-excitation}
\end{equation}
If the width depends only on the source population, so that
$\Sigma_{ij}=\sigma_j^2$, this reduces to
\begin{equation}
 \max_{j\in\mathcal E}\sigma_j
 <
 \min_{j\in\mathcal I}\sigma_j.
 \label{eq:K-source-narrower-excitation}
\end{equation}
For $K>2$, this width ordering is not sufficient by itself.  The off-diagonal
feedback loops affect the critical eigenvectors and eigenvalues, so there is
no universal entrywise strength inequality.  The exact strength condition is
that $\mathsf W(p_c)$ and $\Sigma$ produce
\eqref{eq:K-Hopf-factorization} and
\eqref{eq:K-noncritical-mode-stability} through the matrix
\eqref{eq:K-linear-matrix}.

For $K=2$, the characteristic polynomial is
\begin{equation}
 P(\lambda)=\lambda^2+\gamma_1\lambda+\gamma_2,
 \label{eq:K-two-polynomial}
\end{equation}
and the Hopf conditions reduce to
\begin{equation}
 \gamma_1(k_0,p_c)=0,
 \qquad
 \gamma_2(k_0,p_c)>0.
 \label{eq:K-two-reduction}
\end{equation}
For $K=3$, the characteristic polynomial is
\begin{equation}
 P(\lambda)=\lambda^3+\gamma_1\lambda^2
 +\gamma_2\lambda+\gamma_3,
 \label{eq:K-three-polynomial}
\end{equation}
and the Hopf conditions reduce to
\begin{equation}
 \gamma_1(k_0,p_c)>0,
 \qquad
 \gamma_3(k_0,p_c)>0,
 \qquad
 \gamma_1(k_0,p_c)\gamma_2(k_0,p_c)
 =\gamma_3(k_0,p_c).
 \label{eq:K-three-reduction}
\end{equation}
Thus \eqref{eq:K-two-reduction} gives $A(k_0)=0$, $C(k_0)>0$, while
\eqref{eq:K-three-reduction} gives $A(k_0)>0$, $C(k_0)>0$, and
$A(k_0)B(k_0)=C(k_0)$.
